We derived a formula to extract the expected number of events for the
non-resonant channel needed pre-unblinding to test the signal extraction
strategy:

\begin{align}
n^{\textrm obs}_{\textrm Non} &= 
n^{\textrm obs}_{\textrm Res}
\frac{\mathcal{B}(B\to K^*ee)}{\mathcal{B}(B\to K^*J/\psi)\mathcal{B}(J/\psi\to ee)}
\frac{n_{{\textrm Non},B+\bar B}^{\textrm MC.sel}
\frac{\epsilon^{\textrm gen}_{{\textrm Non},B}\epsilon^{\textrm gen}_{{\textrm Non},\bar{B}}}{n_{{\textrm Non},B}^{\textrm AOD}\epsilon^{\textrm gen}_{{\textrm Non},\bar{B}} + n_{{\textrm Non},\bar{B}}^{\textrm AOD}\epsilon^{\textrm gen}_{{\textrm Non},B}}}{
n_{{\textrm Res},B+\bar B}^{\textrm MC.sel}
\frac{\epsilon^{\textrm gen}_{{\textrm Res},B}\epsilon^{\textrm gen}_{{\textrm Res},\bar{B}}}{n_{{\textrm Res},B}^{\textrm AOD}\epsilon^{\textrm gen}_{{\textrm Res},\bar{B}} + n_{{\textrm Res},\bar{B}}^{\textrm AOD}\epsilon^{\textrm gen}_{{\textrm Res},B}}}
\end{align}

We then define the various part of this formula based on the nature of
the information used. First we consider the external inputs consisting
in the branching ratios of the signal and the resonance channel.

\begin{align}
{\textrm E}(\mathcal{B})_{\ell\ell} &= 
\frac{\mathcal{B}(B\to K^*\ell\ell)}{\mathcal{B}(B\to K^*J/\psi)\mathcal{B}(J/\psi\to \ell\ell)}
\end{align}

This ratio is calculated using the PDG values:

\begin{align}
\mathcal{B}(B\to K^* J/\psi) &= (1.27\pm 0.05)\times 10^{-3}\\
\mathcal{B}(B\to K^* ee) &= (1.04 \pm {0.18})\times 10^{-6}\\
\mathcal{B}(B\to K^*\mu\mu)&= (9.4\pm0.5 )\times 10^{-7}\\
\mathcal{B}(J/\psi \to ee) &= (5.971\pm 0.032) \times 10^{-2}\\
\mathcal{B}(J/\psi \to \mu\mu) &= (5.961\pm 0.033)\times 10^{-2}
\end{align}

and we obtain:

\begin{align}
{\textrm E}(\mathcal{B}_{ee}) = 0.0137 \pm 0.0024 (17.5\%)\\
{\textrm E}(\mathcal{B}_{\mu\mu}) = 0.0124 \pm 0.0008 (6.5\%)
\end{align}

Secondly we consider the efficiencies related to the generation
filters:

\begin{align}
  P({\textrm Non})_{\ell\ell} = \frac{\epsilon^{\textrm gen}_{{\textrm Non},B}\epsilon^{\textrm gen}_{{\textrm Non},\bar{B}}}{n_{{\textrm Non},B}^{\textrm AOD}\epsilon^{\textrm gen}_{{\textrm Non},\bar{B}} + n_{{\textrm Non},\bar{B}}^{\textrm AOD}\epsilon^{\textrm gen}_{{\textrm Non},B}} =
  \frac{1}{\frac{n_{{\textrm Non},B}^{\textrm AOD}}{\epsilon^{\textrm gen}_{{\textrm Non},B}}+\frac{n_{{\textrm Non},\bar{B}}^{\textrm AOD}}{\epsilon^{\textrm gen}_{{\textrm Non},\bar{B}}}}
\end{align}

The numbers needed to calculate the ratios above for the electron
and the muon final states are in Table~\ref{tab:generation}.

\begin{table}[!tbp]
\begin{center}
\smallskip
\begin{tabular}{|l|c|c|} 
\hline
Inputs to the P ratios: & & \\
\hline
Non-resonant channel & & \\
\hline
\hline
generator efficiencies & B & $\overline{B}$ \\
\hline
electron channel & $0.05335$ & $0.05289$ \\
muon channel & $0.06442$ & $0.06403$ \\
\hline
Number of events & B & $\overline{B}$ \\
\hline
electron channel & $6441500$ & $4976000$ \\
muon channel & $1498000$ & $1499000$ \\
\hline
\hline
Resonant channel & & \\
\hline
\hline
generator efficiencies & B & $\overline{B}$ \\
\hline
electron channel & $0.05850$ & $0.05802$ \\
muon channel & $0.06137$ & $0.06084$ \\
\hline
Number of events & B & $\overline{B}$ \\
\hline
electron channel & $447500$ & $449500$ \\
muon channel & $436000$ & $449500$ \\
\hline
\end{tabular}
\caption{Numbers of generated events and generation efficiencies for
the non-resonant and resonant signals.}
\label{tab:generation}
\end{center}
\end{table}

%% to get to the original unfiltered number of generated events
%% from the generation efficiency and the number of actually generated events.
%6441500/0.05335
%120740394
%sqrt(0.05335*(1-0.05335)/120740394)
%.00002045199308136006
%0.00002/0.05335
%.00037488284910965323
%% the error is very small so we consider these numbers without errors

The resulting ratios from the generation efficiencies are as follow:

\begin{align}
{\textrm P}(\textrm Non)_{ee} = 0.466 \times 10^{-8}\\
{\textrm P}(\textrm Res)_{ee} = 6.495 \times 10^{-8}\\
{\textrm P}(\textrm Non)_{\mu\mu} = 2.143 \times 10^{-8}\\
{\textrm P}(\textrm Res)_{\mu\mu} = 6.900 \times 10^{-8}
\end{align}

We now can rewrite the main formula in a simplify way:
\begin{align}
n^{\textrm obs}_{\textrm Non} &= 
n^{\textrm obs}_{\textrm Res}
{\textrm E}(\mathcal{B})
\frac{n_{{\textrm Non},B+\bar B}^{\textrm MC.sel}}{n_{{\textrm Res},B+\bar B}^{\textrm MC.sel}}
\frac{P({\textrm Non})}{P({\textrm Res})}
\end{align}
where we are left to input the number of events after
the final selection.

%\begin{align}
%  n_{{\textrm Non},B+\bar B}^{\textrm MC.sel}\mbox{ }_{ee} = 34867\\
%  n_{{\textrm Res},B+\bar B}^{\textrm MC.sel}\mbox{ }_{ee} = 31576\\
%  n_{{\textrm Non},B+\bar B}^{\textrm MC.sel}\mbox{ }_{\mu\mu} = \\
%  n_{{\textrm Res},B+\bar B}^{\textrm MC.sel}\mbox{ }_{\mu\mu} = 
%\end{align}

Assuming a GNN working point for the electron final states
and a BDT working point for the muon final states
(give a couple of numbers to characterise the WP:
signal efficiency and background efficiency?),
we obtain the following number of events:
\begin{align}
  n_{{\textrm Non},{B+\bar B}}^{\textrm MC,sel}(ee) &= 34867\\
  n_{{\textrm Res},{B+\bar B}}^{\textrm MC,sel}(ee) &= 31576\\
  n_{{\textrm Non},B+\bar{B}}^{\textrm MC, sel}(ee,weighted) &= 838 \pm 29\\
  n_{{\textrm Res},B+\bar{B}}^{\textrm MC, sel}(ee,weighted) &= 841 \pm 28\\
  n_{{\textrm Non},{B+\bar B}}^{\textrm MC,sel}(\mu\mu) &= 4.3 \times 10^4\\
  n_{{\textrm Res},{B+\bar B}}^{\textrm MC,sel}(\mu\mu) &= 7.9 \times 10^4\\
  n_{{\textrm Non},B+\bar{B}}^{\textrm MC, sel}(\mu\mu,weighted) &=  \pm \\
  n_{{\textrm Res},B+\bar{B}}^{\textrm MC, sel}(\mu\mu,weighted) &=  \pm 
\end{align}
where we include the raw numbers of events passing in the MC and
also the weighted numbers of events. We should use the weighted numbers,
when available.

From the resonant fit to the B invariant mass we obtain the follow
number of events:
\begin{align}
  n^{\textrm obs}_{\textrm Res}(ee) &= 12359 \pm 217 \\
  n^{\textrm obs}_{\textrm Res}(\mu\mu) &= 9.4 \times 10^4 \pm ??\\
\end{align}

Applying the formula we then obtain the following number of expected
events for the non-resonant channel:
\begin{align}
  n^{\textrm exp}_{\textrm Non}(ee) &= 24 \pm 10 (40\%)\\
  n^{\textrm exp}_{\textrm Non}(\mu\mu) &= 210 (15\%)\\
\end{align}

We can now proceed to estimate the factors contributing to the
extraction of R(K*):

\begin{align}
  R(K^\ast) &= \frac{\frac{\mathcal{B}(B\to K^\ast\mu\mu)}{\mathcal{B}(B\to K^\ast J/\psi)}}
  {\frac{\mathcal{B}(B\to K^\ast ee)}{\mathcal{B}(B\to K^\ast J/\psi)}}
\end{align}

For electrons:
\begin{align}
  {\textrm{ratio of gen efficiencies non/res}} &=  0.0535/0.0585 = 0.92  \\
  {\textrm{ratio of preselection efficiencies non/res}} &= 0.032/0.155 = 0.21 \\
  {\textrm{ratio of trigger efficiencies non/res}} &=  0.0295/0.0317 = 0.93\\
  {\textrm{ratio of GNN efficiencies non/res}} &=   \\
\end{align}

The lower preselection efficiency for non-resonant events versus resonant is
physically motivated by both the $q^2$ selection and the $Delta R(ee)$ cut where in the
non-resonant events the $0.1$ requirement cuts away the photon pole at zero.

For muons:
\begin{align}
  {\textrm{ratio of gen efficiencies non/res}} &=   0.0644/0.0614 = 1.05  \\
  {\textrm{ratio of preselection efficiencies non/res}} &=  \\
  {\textrm{ratio of trigger efficiencies non/res}} &=  \\
  {\textrm{ratio of GNN efficiencies non/res}} &=  \pm \\
\end{align}

Taking the numbers into account the back-of-the-envelope consideration is that
the uncertainty in the R(K*) value is dominated by the uncertainty on the
non-resonant yield in the electron channel (currently at 44\%). This would be
comparable with CMS which has a 46\% error in RK.

We consider also the measurement of the BR of the non-resonant channel in the
muon final state. Here the error will be also dominated by the statistical
uncertainty on the yield. Considering a number of events at 200, we get a 15\%
uncertainty. However this is only a subset of the 2018 dataset so we aim at being
comparable with CMS which has a 5\% error in the low-$q^2$ bin.


